(1) 沿环的正则角动量算符是 \(-i\hbar\partial_\theta\)，机械动量需扣去 \(qRA_\theta=q\Phi/(2\pi)\)。因此 \(\boxed{H=\frac1{2mR^2}(-i\hbar\partial_\theta-q\Phi/(2\pi))^2}\)，它只依赖包围的磁通，不要求粒子运动位置有非零局域磁场。

(2) 单值波函数 \(\psi_n(\theta)=e^{in\theta}/\sqrt{2\pi}\)、\(n\in\mathbb Z\)。记 \(\varphi=q\Phi/h\)，则 \(\boxed{E_n(\Phi)=\frac{\hbar^2}{2mR^2}(n-\varphi)^2}\)。磁通改变一个量子 \(h/|q|\) 可用整数 \(n\) 重标记全部能谱，所以集合和热力学状态具有磁通周期。

(3) 每个本征态的持久电流为 \(\boxed{I_n=-\frac{\partial E_n}{\partial\Phi}=\frac{q\hbar}{2\pi mR^2}(n-\varphi)}\)。机械角速度是 \(\dot\theta=\hbar(n-\varphi)/(mR^2)\)，一粒电荷每圆周的电流是 \(q\dot\theta/(2\pi)\)，与热力学导数完全一致。

(4) 基态选择最接近 \(\varphi\) 的整数 \(n_*\)，因此 \(E_G=\hbar^2(n_*-\varphi)^2/(2mR^2)\) 是周期性分段抛物线，\(I_G\) 为分段线性锯齿函数。在 \(\varphi\in\mathbb Z+1/2\) 时相邻两个整数能级相同，理想零温基态简并并产生电流跳变；无序势、破缺旋转对称或有限温混合可能平滑这些非解析行为。
